% ============================================================
% PROYECTO INTEGRADOR: SMS DEL TALLER DE AVIÓNICA
% ============================================================
\clearpage
\section*{PROYECTO INTEGRADOR}
\subsection*{Sistema de Gestión de la Seguridad Operacional del taller de aviónica}

\begin{recuadro}
\textbf{Consigna general:} diseñar e implementar, durante el segundo cuatrimestre, un
\textbf{SMS simplificado} para el taller de aviónica de la escuela, siguiendo los cuatro
componentes del Anexo 19 de la OACI, adaptado al tamaño y a la complejidad del taller.

\textbf{Modalidad:} trabajo en equipos con roles. \textbf{Entrega final:} Manual SMS del taller,
registro de peligros y matriz de riesgos, sistema de reportes funcionando, tablero de indicadores
y presentación oral ante la dirección de la escuela.
\end{recuadro}

% ------------------------------------------------------------
\section*{1. ORGANIZACIÓN}
\begin{center}
\small
\begin{tabular}{|p{0.28\linewidth}|p{0.62\linewidth}|}
\hline
\textbf{Equipo} & \textbf{Responsabilidad} \\ \hline
Política y documentación & Redactar la política de seguridad, definir roles y armar el Manual
SMS. \\ \hline
Gestión de riesgos & Identificar peligros, evaluarlos en la matriz y proponer mitigaciones. \\
\hline
Reportes e investigación & Diseñar el sistema de reportes voluntarios e investigar los sucesos
con el método MEDA. \\ \hline
Aseguramiento & Definir indicadores, metas, auditorías internas y gestión del cambio. \\ \hline
Promoción & Capacitar a otros cursos, publicar un boletín de seguridad mensual y la cartelería. \\
\hline
\end{tabular}
\end{center}
El \textbf{comité de seguridad operacional} reúne a un representante de cada equipo y al docente
(que actúa como responsable del SMS), y se reúne cada quince días.

% ------------------------------------------------------------
\section*{2. ETAPAS}

\subsubsection*{Etapa 1. Política y objetivos (Componente 1)}
\begin{itemize}
  \item Redactar la \textbf{política de seguridad operacional} del taller (una página): compromiso,
        cultura justa, obligación de reportar, objetivos. Firmada por la dirección o el jefe de
        taller.
  \item Definir responsabilidades: quién es responsable del SMS, qué hace cada estudiante y cada
        docente.
  \item Elaborar el \textbf{plan de respuesta ante emergencias} del taller: incendio,
        electrocución, accidente con herramientas, evacuación; números de emergencia, ubicación de
        matafuegos, botiquín y DEA.
\end{itemize}

\subsubsection*{Etapa 2. Gestión de riesgos (Componente 2)}
\begin{itemize}
  \item Recorrer el taller y armar el \textbf{registro de peligros} (mínimo 15).
  \item Evaluar cada uno en la \textbf{matriz de riesgos} y proponer medidas con la jerarquía de
        controles.
\end{itemize}

\begin{center}
\footnotesize
\setlength{\tabcolsep}{3pt}
\newcolumntype{L}[1]{>{\raggedright\arraybackslash}p{#1\linewidth}}
\begin{tabular}{|c|L{0.21}|L{0.17}|c|L{0.21}|L{0.08}|L{0.11}|}
\hline
\textbf{N°} & \textbf{Peligro} & \textbf{Consecuencia} & \textbf{Riesgo} & \textbf{Mitigación} &
\textbf{Riesgo resid.} & \textbf{Respons.} \\ \hline
1 & Fuente de 28 V DC sin protección en el banco & Descarga, quemadura, daño a equipos & 3C &
Tapas protectoras, bloqueo y etiquetado & 2D & Equipo 2 \\ \hline
2 & Herramientas sin control al terminar la clase & Herramienta olvidada dentro de un equipo
(FOD) & 4C & Tablero de sombra y conteo final & 2D & Pañol \\ \hline
3 & & & & & & \\ \hline
4 & & & & & & \\ \hline
\end{tabular}
\end{center}

\subsubsection*{Etapa 3. Aseguramiento (Componente 3)}
\begin{itemize}
  \item Definir al menos \textbf{cuatro indicadores} (SPI) con su meta, por ejemplo: reportes
        voluntarios por mes; herramientas faltantes al cierre de la clase; equipos devueltos con
        falla al pañol; hallazgos de la auditoría interna resueltos.
  \item Hacer una \textbf{auditoría interna} del taller con una lista de verificación.
  \item Aplicar la \textbf{gestión del cambio} a una novedad real (un equipo nuevo, un cambio de
        aula, un nuevo procedimiento): ¿qué peligros nuevos aparecen?
\end{itemize}

\subsubsection*{Etapa 4. Promoción (Componente 4)}
\begin{itemize}
  \item Dar una \textbf{capacitación} de 40 minutos a un curso menor sobre la docena sucia y el
        sistema de reportes.
  \item Publicar un \textbf{boletín de seguridad} mensual con los reportes recibidos (sin nombres),
        las lecciones aprendidas y los indicadores.
\end{itemize}

% ------------------------------------------------------------
\section*{3. FORMULARIO DE REPORTE VOLUNTARIO}
{\small
\renewcommand{\arraystretch}{1.5}
\begin{tabular}{|p{0.955\linewidth}|}
\hline
\textbf{REPORTE VOLUNTARIO DE SEGURIDAD OPERACIONAL -- TALLER DE AVIÓNICA} \hfill N°: \makebox[2cm]{\dotfill} \\ \hline
Fecha y hora del suceso: \dotfill \quad Lugar: \makebox[4.5cm]{\dotfill} \\ \hline
Tipo: \casilla\ Peligro \quad \casilla\ Error o casi accidente \quad \casilla\ Incidente \quad
\casilla\ Sugerencia \\ \hline
¿Qué pasó o qué observó? \\[1.6cm] \hline
¿Por qué cree que pasó? \\[1.2cm] \hline
¿Qué propone para que no vuelva a pasar? \\[1.2cm] \hline
Nombre (optativo): \dotfill \quad \casilla\ Deseo recibir respuesta \\ \hline
\textit{Este reporte es confidencial. Su objetivo es aprender y prevenir, no sancionar.} \\ \hline
\end{tabular}
}

% ------------------------------------------------------------
\section*{4. ENTREGABLES Y EVALUACIÓN}
\begin{center}
\small
\begin{tabular}{|p{0.66\linewidth}|c|}
\hline
\textbf{Entregable} & \textbf{Peso} \\ \hline
Política de seguridad, roles y plan de respuesta ante emergencias & 15\% \\ \hline
Registro de peligros y matriz de riesgos con mitigaciones & 25\% \\ \hline
Sistema de reportes funcionando e investigación MEDA de al menos un suceso & 20\% \\ \hline
Indicadores, auditoría interna y gestión del cambio & 15\% \\ \hline
Capacitación y boletín de seguridad & 10\% \\ \hline
Manual SMS completo y presentación oral & 15\% \\ \hline
\end{tabular}
\end{center}

\begin{recuadro}
\textbf{Para el CV:} al terminar, cada estudiante puede incluir en su currículum la experiencia
de \textit{``Diseño e implementación de un Sistema de Gestión de la Seguridad Operacional (SMS)
según el Anexo 19 de la OACI en un taller de aviónica''}, junto con las tareas concretas que
realizó en su equipo.
\end{recuadro}
